\chapter{참고문헌}
\begingroup\renewcommand{\chapter}[2]{}\begin{thebibliography}{99}
\bibitem{b1} I. Aghabali, J. Bauman, P. J. Kollmeyer, Y. Wang, B. Bilgin, A. Emadi, “800-V Electric Vehicle Powertrains: Review and Analysis of Benefits, Challenges, and Future Trends,” IEEE Transactions on Transportation Electrification, vol. 7, no. 3, pp. 927–948, 2021. DOI: 10.1109/TTE.2020.3044938.
\bibitem{b2} B. Zhao, Q. Song, W. Liu, Y. Sun, “Overview of Dual-Active-Bridge Isolated Bidirectional DC–DC Converter for High-Frequency-Link Power-Conversion System,” IEEE Transactions on Power Electronics, vol. 29, no. 8, pp. 4091–4106, 2014. DOI: 10.1109/TPEL.2013.2289913.
\bibitem{b3} Z. Ni, Y. Li, C. Liu, M. Wei, D. Cao, “A 100-kW SiC Switched Tank Converter for Transportation Electrification,” IEEE Transactions on Power Electronics, vol. 35, no. 6, pp. 5770–5784, 2020. DOI: 10.1109/TPEL.2019.2954801.
\bibitem{b4} Y. He, S. Jiang, C. Nan, “Switched Tank Converter Based Partial Power Architecture for Voltage Regulation Applications,” IEEE APEC, pp. 91–97, 2018. DOI: 10.1109/APEC.2018.8340993.
\bibitem{b5} Z. Ni, Y. Li, J. Johnson, M. Wei, C. Liu, X. Lyu, D. Cao, “A GaN Switched Tank Converter with Partial Power Voltage Regulation for Electric Vehicle Applications,” IEEE APEC, pp. 1682–1689, 2019. DOI: 10.1109/APEC.2019.8722301.
\bibitem{b6} P. Wang, M. Chen, “Analysis and Design of Series Voltage Compensator for Differential Power Processing,” IEEE Journal of Emerging and Selected Topics in Power Electronics, vol. 10, no. 6, pp. 7890–7903, 2022. DOI: 10.1109/JESTPE.2021.3116091.
\bibitem{b7} S. Bibian, H. Jin, “Time Delay Compensation of Digital Control for DC Switchmode Power Supplies Using Prediction Techniques,” IEEE Transactions on Power Electronics, vol. 15, no. 5, pp. 835–842, 2000. DOI: 10.1109/63.867672.
\bibitem{b8} I. Kolmanovsky, E. Garone, S. Di Cairano, “Reference and Command Governors: A Tutorial on Their Theory and Automotive Applications,” American Control Conference, pp. 226–241, 2014. DOI: 10.1109/ACC.2014.6859176.
\bibitem{b9} ISO 21498-1:2021, “Electrically propelled road vehicles — Electrical specifications and tests for voltage class B systems and components — Part 1: Voltage sub-classes and characteristics,” Table 1.
\bibitem{b10} ISO 21780:2020, “Road vehicles — Supply voltage of 48 V — Electrical requirements and tests,” Clause 5.
\bibitem{b11} H. Chen, H. Kim, R. Erickson, D. Maksimovi{\fallbackfont ć}, “Electrified Automotive Powertrain Architecture Using Composite DC–DC Converters,” IEEE Transactions on Power Electronics, vol. 32, no. 1, pp. 98–116, 2017.
\bibitem{b12} Wolfspeed, “C3M0016120K Silicon Carbide Power MOSFET, 1200 V, 16 mΩ,” datasheet.
\bibitem{b13} onsemi (UnitedSiC), “UJ4SC075006K4S 750 V, 5.9 mΩ G4 SiC FET,” datasheet.
\bibitem{b14} GaN Systems (Infineon), “GS66516T Top-side cooled 650 V E-mode GaN transistor,” datasheet Rev. 210727. 2026-10-28 열람.
\bibitem{b15} Efficient Power Conversion, “EPC2305 – Enhancement Mode Power Transistor (150 V, 2.2 mΩ),” datasheet(I\textsubscript{D,pulse} 329 A at 25 °C·300 µs, V\textsubscript{SD} 1.4 V). 2026-10-28 열람.
\bibitem{b16} IEEE Std 2030.1.1-2015, “IEEE Standard Technical Specifications of a DC Quick Charger for Use with Electric Vehicles.” 시험 항목 요약: IEEE SA, “Test Cases for DC Quick Charger for Use with Electric Vehicles”(항목 4.5.2.5.4 a)·c): 충전 허가 신호가 꺼지거나 커넥터 잠금 결함을 검출한 뒤 30 ms 안에 출력전류 5 A 이하, 4.7.3.4·4.5.2.1.4 c): 5 A 이하까지 접촉기 회로 유지, 4.5.1.15.4: 접촉기 차단 시 커넥터 전압 ≤ 600 V(서지 제외). 모두 부속서 A(CHAdeMO 사양)를 인용하며, 대상 충전기 출력은 150–500 V).
\bibitem{b17} J.-H. Jung, H.-S. Kim, M.-H. Ryu, J.-W. Baek, “Design Methodology of Bidirectional CLLC Resonant Converter for High-Frequency Isolation of DC Distribution Systems,” IEEE Transactions on Power Electronics, vol. 28, no. 4, pp. 1741–1755, 2013. DOI: 10.1109/TPEL.2012.2213346.
\bibitem{b18} N. Zanatta, T. Caldognetto, D. Biadene, G. Spiazzi, P. Mattavelli, “Design and Implementation of a Two-Stage Resonant Converter for Wide Output Range Operation,” IEEE Transactions on Industry Applications, 2022. DOI: 10.1109/TIA.2022.3199195.
\bibitem{b19} C.-E. Kim, “Minimization of Abnormal Output Voltage Rising for LLC Resonant Converter at Very Light Load,” IEEE Transactions on Industrial Electronics, vol. 67, no. 12, pp. 10295–10303, 2020. DOI: 10.1109/TIE.2019.2960730.
\bibitem{b20} R. D. Middlebrook, “Input Filter Considerations in Design and Application of Switching Regulators,” Proc. IEEE Industry Applications Society Annual Meeting, 1976.
\bibitem{b21} A. Emadi, A. Khaligh, C. H. Rivetta, G. A. Williamson, “Constant Power Loads and Negative Impedance Instability in Automotive Systems: Definition, Modeling, Stability, and Control of Power Electronic Converters and Motor Drives,” IEEE Transactions on Vehicular Technology, vol. 55, no. 4, pp. 1112–1125, 2006. DOI: 10.1109/TVT.2006.877483.
\bibitem{b22} K. Kesarwani, R. Sangwan, J. T. Stauth, “Resonant-Switched Capacitor Converters for Chip-Scale Power Delivery: Design and Implementation,” IEEE Transactions on Power Electronics, vol. 30, no. 12, pp. 6966–6977, 2015.
\bibitem{b23} Y. Ahn, I. Jeon, J. Roh, “A Multiphase Buck Converter With a Rotating Phase-Shedding Scheme For Efficient Light-Load Control,” IEEE Journal of Solid-State Circuits, vol. 49, no. 11, pp. 2673–2683, 2014. DOI: 10.1109/JSSC.2014.2360400.
\bibitem{b24} Hyundai Motor Group, E-GMP 플랫폼 소개(제조사 공개 자료), 2020.
\bibitem{b25} Lucid Group, “Lucid Gravity: Combining 450 miles EPA Range, 400 kW Fast Charging, and Access to 20,000+ Superchargers,” 보도자료, 2025.
\bibitem{b26} Preh GmbH, “High-Voltage Booster — Efficient charging of 800V vehicles at 400V charging stations,” 제품 소개(제조사 공개 자료).
\bibitem{b27} G. Qiu, V. Khadkikar, M. R. Mohammed, B. Zahawi, “Power Transfer from 400V Charging Piles to 800V Electric Vehicles Using Motor Winding and Inverter,” Proc. IEEE PEDES, 2024. DOI: 10.1109/PEDES61459.2024.10961770.
\bibitem{b28} M. Eull, L. Zhou, M. Jahnes, M. Preindl, “Bidirectional Nonisolated Fast Charger Integrated in the Electric Vehicle Traction Drivetrain,” IEEE Transactions on Transportation Electrification, vol. 8, no. 1, pp. 180–195, 2022. DOI: 10.1109/TTE.2021.3124936.
\bibitem{b29} Porsche AG, “The charging process: Quick, comfortable, intelligent and universal,” Porsche Newsroom, 2019.
\bibitem{b30} M. Shen, F. Guo (NIO USA), “Systems and methods for regulating charging of electric vehicles,” US Patent 10,363,828 B1 / US 10,632,863 B2.
\bibitem{b31} Karma Automotive, “Combined converter circuit,” US Patent 11,235,676 B2 / EP 3987646.
\bibitem{b32} N. Perera, A. Jafari, R. Soleimanzadeh, N. Bollier, S. G. Abeyratne, E. Matioli, “Hard-Switching Losses in Power FETs: The Role of Output Capacitance,” IEEE Transactions on Power Electronics, vol. 37, no. 7, pp. 7604–7616, 2022. DOI: 10.1109/TPEL.2021.3130831.
\bibitem{b33} H. B. Sambo, R. A. Abramson, S. Krishnan, R. C. N. Pilawa-Podgurski, “An Adaptive Zero Current Switching Control Technique for Multi-Resonant Switched-Capacitor Converters,” Proc. IEEE APEC, 2025.
\bibitem{b34} Wolfspeed, “C4MS008075U2 / C4MS025075K1 Silicon Carbide Power MOSFET, 750 V,” datasheets Rev. 1, 2026; onsemi, “UG4SC075011K4S 750 V SiC FET,” datasheet; GaN Systems, “GS66516T,” datasheet Rev. 210727. 2026-10-28 열람.
\bibitem{b35} M. Jenks (Vicor), “Adapting 400V and 800V Architectures,” SAE WCX 2023 발표 자료, 2023. 2026-10-27 열람.
\bibitem{b36} S. Tsukimoto (Vicor), “Power modules provide high-efficiency conversion between 400V and 800V systems for electric vehicles,” Vicor 기술 기사, 작성일 미표기. 2026-10-27 열람.
\bibitem{b37} Vicor, “Solving 400V/800V EV fast-charging compatibility”(NBM9280 소개), 제조사 자료, 작성일 미표기. 2026-10-27 열람.
\bibitem{b38} Vicor, “Supporting legacy accessories in new 800V architectures — Case study: Powering 400V accessory loads,” Rev 1.2, 2025. 2026-10-27 열람.
\bibitem{b39} M. Shen, F. Guo (NIO USA), “Devices, systems, and methods for charging electric vehicles,” US Patent 10,924,008 B2, 2021.
\bibitem{b40} M. Uno, Y. Fujii, T. Sanada, “ZVS Resonant Switched Capacitor Converter with 1.5x or 2.0x Step-Up Conversion Ratio for High-Voltage EV Batteries and Powertrains,” IEEJ Journal of Industry Applications, vol. 14, no. 1, pp. 103–110, 2025. DOI: 10.1541/ieejjia.24006436.
\bibitem{b41} C. T. Lieberman, P. Saggere, H. Mhiesan, C. Ballard (Ford Global Technologies), “Electrified vehicle with synchronous switching converter supporting exportable power and DC fast charging,” US Patent 12,179,618 B1, 2024.
\bibitem{b42} D. Li, L. Hao, C. J. Lee, S. Gopalakrishnan, J. Zhao (GM Global Technology Operations), “High voltage switching for an electric vehicle,” US Patent 11,865,932 B2, 2024.
\bibitem{b43} C. S. Namuduri, R. Prasad (GM Global Technology Operations), “High voltage switching for charging an electric vehicle,” US Patent 11,840,153 B1, 2023.
\bibitem{b44} A. Lasson, R. Eriksson (Volvo Car), “Charging system and method for charging a battery of an electric vehicle,” EP 4 001 001 A1, 2022.
\bibitem{b45} W. Y. Lee et al. (Hyundai Motor, Kia Motors), “Electric vehicle,” US Patent 10,358,041 B2, 2019.
\bibitem{b46} L. Zhu, H. Bai, A. Brown, “Model and Control of a Current-Fed Dual Active Bridge Based Ultrawide-Voltage-Range Auxiliary Power Module for 400 V/800 V Electric Vehicles,” IEEE Transactions on Power Electronics, vol. 39, no. 3, pp. 3263–3276, 2024. DOI: 10.1109/TPEL.2023.3337712.
\bibitem{b47} GM Global Technology Operations, “Energy transfer among mismatched battery systems,” US Patent 12,325,327 B2, 2025.
\bibitem{b48} A. Blinov, A. Chub, O. Korkh, D. Vinnikov, I. Galkins, D. Peftitsis (Tallinn University of Technology), “System and method for a partial power transfer between two DC sources,” US Patent 12,301,121.
\bibitem{b49} N. Hassanpour, A. Blinov, A. Chub, D. Vinnikov, “Soft start and protection of bidirectional buck-boost partial power converter,” Proc. IEEE SGRE, pp. 1–6, 2022. DOI: 10.1109/SGRE53517.2022.9774133.
\bibitem{b50} N. Hassanpour, A. Chub, A. Blinov, D. Vinnikov, “Soft-Switching Bidirectional Step-Up/Down Partial Power Converter With Reduced Components Stress,” IEEE Transactions on Power Electronics, vol. 38, no. 11, pp. 14166–14177, 2023. DOI: 10.1109/TPEL.2023.3289061.
\bibitem{b51} A. Siddiquee, A. Uzum, S. Hasan, Y. Sozer, M. J. Kisacikoglu, “Traction Inverter Integrated On-Board DC Fast Charging through Partial Power Processing,” Proc. IEEE ECCE, pp. 2134–2139, 2024. DOI: 10.1109/ECCE55643.2024.10861726.
\bibitem{b52} B. Ejaz, R. Zamora, C. Reusser, X. Lin, “A Comprehensive Review of Partial Power Converter Topologies and Control Methods for Fast Electric Vehicle Charging Applications,” Electronics, vol. 14, no. 10, 1928, 2025. DOI: 10.3390/electronics14101928.
\bibitem{b53} A. Babazadeh, R. Renganathan, D. Clavette, C. Rainer (Infineon Technologies Austria), “Method and apparatus for zero-current switching control in switched-capacitor converters,” US Patent 10,122,256 B1, 2018.
\bibitem{b54} Y. Li, X. Lyu, Z. Ni, D. Cao, C. Nan, S. Jiang, “Adaptive On-Time Control for High Efficiency Switched-Tank Converter,” Proc. IEEE WiPDA Asia, pp. 169–175, 2018. DOI: 10.1109/WiPDAAsia.2018.8734623.
\bibitem{b55} R. L. Valley, I. Cohen (Texas Instruments), “System and method for phase management in a multiphase switching power supply,” US Patent 8,120,334 B2, 2012.
\bibitem{b56} Infineon Technologies, “FF2MR12KM1P 62 mm C-Series CoolSiC Trench MOSFET half-bridge module,” datasheet V2.0(Preliminary Data), 2020-05-11. 2026-10-29 열람.
\bibitem{b57} Wolfspeed, “CAB530M12BM3 1200 V, 530 A all-SiC half-bridge module,” data sheet Rev. 5; “CAB450M12XM3 1200 V, 450 A all-SiC half-bridge module,” data sheet Rev. A, 2019. 2026-10-28·29 열람.
\bibitem{b58} Ohmite, “800 and 1000 Series Tubular High Energy Bulk Ceramic Resistors,” rev 7/25-1; Vishay Sfernice, “LTO 150H,” datasheet 및 “Thick Film Power Resistors: D2TO-H and LTO-H High Energy Series,” Application Note 51114, 2025. 2026-10-28·29 열람.
\bibitem{b59} Fischer \& Tausche Capacitors, “Film Capacitors — Electrical Data for Series CASTOR CAP,” 개정일 미표기; TDK Electronics, “B32774/B32776/B32778 Film Capacitors for DC Link,” datasheet. 2026-10-28·29 열람.
\bibitem{b60} TE Connectivity, “High Voltage DC Contactor ECK100B Series,” datasheet, 2024. 2026-10-28 열람.
\end{thebibliography}\endgroup
